%% ========================================================================
%% chapter04-directory-structure.tex
%% Directory Structure and File Operations
%% ========================================================================

\chapter{Directory Structure and File Operations}
\label{ch:filesystem-ops}

\chapterhero{orange}{Explore the Linux file system structure, understand file types and links, then master searching and archiving for daily work.}{\faFolderOpen\quad FHS}{\faLink\quad symbolic link}{\faArchive\quad archive}

The system you installed in Chapter~1 already contains thousands of files spread across hundreds of directories. This is not disorder. It is a pattern standardized by the Linux community. An administrator who is familiar with Ubuntu can quickly find configuration files on a Debian machine they have never touched, because both follow the same layout rules.

The Linux \term{filesystem} is not a pile of files placed randomly. Every directory has a defined role. Application configuration files live in \dir{/etc/}, system activity records are stored in \dir{/var/log/}, and programs executable by all users are in \dir{/bin/}. Once you know these rules, finding files in Linux becomes fast and systematic, not guesswork.

In addition to directory layout, Linux distinguishes files by type: regular files, directories, symbolic links, device files, and more. These differences are not just labels. They determine how the system treats each object. Understanding them helps you diagnose access problems, design tidy project folder structures, and work with Linux systems more confidently.

\textbf{What You'll Learn}

By the end of this chapter, you will be able to:

\begin{enumerate}
    \item explain the function of major directories in the Linux file system hierarchy;
    \item distinguish Linux file types based on the first character in \cmd{ls -l} output;
    \item explain the difference between hard links and symbolic links and create them with the correct commands;
    \item search for files by name, type, size, and modification time using \cmd{find} and \cmd{locate};
    \item create, inspect, and extract compressed archives with \cmd{tar} and \cmd{gzip}.
\end{enumerate}

\begin{notebox}
All labs in this chapter use the directory \dir{~/lab-os/chapter04-files/}.
Create this directory before starting:
\begin{lstlisting}[language=bash]
mkdir -p ~/lab-os/chapter04-files
cd ~/lab-os/chapter04-files
\end{lstlisting}
\end{notebox}

% ------------------------------------------------------------------------
\section{Filesystem Hierarchy Standard}
\label{sec:fhs}

Linux organizes all files in a single directory tree that starts from the root, \dir{/}. There is no ``drive C'' or ``drive D'' concept as in Windows. All storage devices, partitions, and even network \term{filesystem} instances are mounted into the same tree. The standard that governs this layout is called the \textit{Filesystem Hierarchy Standard} (FHS).

The main Unix principle to remember is: \textit{everything is a file}. This does not apply only to text documents. Hardware such as hard disks is represented as files in \dir{/dev/}. Information about running processes can be read like files in \dir{/proc/}. Understanding this principle opens a new way to see how Linux works.

Here are the major directories you need to know. Linux administrators memorize these naturally through daily use.

\begin{itemize}
    \item \dir{/} : the starting point of the entire hierarchy. All other directories come from here.
    \item \dir{/bin/} : essential programs needed by all users, such as \cmd{ls}, \cmd{cp}, \cmd{mv}, and \cmd{bash}.
    \item \dir{/sbin/} : system administration programs, such as \cmd{fdisk} and \cmd{iptables}. These usually require \cmd{sudo}.
    \item \dir{/etc/} : all system configuration files. They are text-based, so they can be read with an ordinary editor.
    \item \dir{/home/} : each user's personal directory. User \texttt{alice} has \dir{/home/alice/}.
    \item \dir{/root/} : the personal directory for the \texttt{root} user. It is separated from \dir{/home/} for security.
    \item \dir{/var/} : data that keeps changing while the system runs. The most important subdirectory is \dir{/var/log/}, which stores system logs.
    \item \dir{/tmp/} : temporary files. Usually cleared automatically when the system restarts.
    \item \dir{/usr/} : programs and \term{shared library} files for users. \dir{/usr/bin/} contains most applications you install.
    \item \dir{/opt/} : third-party applications installed independently outside the package manager.
    \item \dir{/dev/} : device files. \file{/dev/sda} is the first hard disk, and \file{/dev/tty} is a terminal.
    \item \dir{/proc/} : a virtual filesystem from the kernel. Reading it gives information about processes and current kernel configuration.
    \item \dir{/sys/} : another virtual \term{filesystem} from the kernel, containing hardware and \term{driver} information.
    \item \dir{/mnt/} and \dir{/media/} : mount points for external devices such as USB drives or additional partitions.
\end{itemize}

\subsection*{Lab 4.1: Explore Major Directories}

The goal of this lab is to train your eyes to recognize standard Linux directories and directly find important configuration files.

\begin{enumerate}
    \item Display the root directory contents and notice names you already know:
\begin{lstlisting}[language=bash, caption={Viewing the root filesystem contents}]
ls -lah /
\end{lstlisting}
    \textit{Observe:} Notice the first column of each line. The character \texttt{d} marks a directory, while \texttt{l} marks a symbolic link.

    \item View the three directories administrators visit most often:
\begin{lstlisting}[language=bash, caption={Inspecting /etc, /var/log, and /home}]
ls /etc | head -20
ls /var/log | head -20
ls /home
\end{lstlisting}

    \item Read CPU information directly from the kernel virtual filesystem:
\begin{lstlisting}[language=bash, caption={Reading system information from /proc}]
cat /proc/cpuinfo | head -20
cat /proc/meminfo | head -10
\end{lstlisting}
    \textit{Notice:} Files in \dir{/proc/} feel like ordinary files, but their contents are generated directly by the kernel, not stored on disk.

    \item Identify three important configuration files in \dir{/etc/}:
\begin{lstlisting}[language=bash, caption={Finding important configuration files with grep}]
ls /etc | grep -E "^(hostname|hosts|passwd|group|fstab)$"
cat /etc/hostname
\end{lstlisting}
\end{enumerate}

\begin{challengebox}
Use a combination of \cmd{ls} and \cmd{grep}, which you learned in Chapter~2, to find three configuration files in \dir{/etc/} whose names contain \texttt{ssh}. After finding them, display the first 10 lines of each file using \cmd{head}. Document the commands you used and write a short function for each file.
\end{challengebox}

% ------------------------------------------------------------------------
\section{File Types in Linux}
\label{sec:filetypes}

When you run \cmd{ls -l}, the first character on each line is not decoration. It tells you the file type. Linux recognizes seven file types, and distinguishing them matters because the system behaves differently for each type.

Here are the file type marker characters you will often see:

\begin{itemize}
    \item \texttt{-} : regular file. Text documents, images, and program \term{binary} files all belong to this category.
    \item \texttt{d} : directory. A container that can hold files and other directories.
    \item \texttt{l} : symbolic link. A shortcut that points to another file or directory.
    \item \texttt{b} : block device file. An example is a hard disk partition such as \file{/dev/sda1}.
    \item \texttt{c} : character device file. An example is the terminal \file{/dev/tty}.
    \item \texttt{p} : named pipe. Used for inter-process communication.
    \item \texttt{s} : socket. Used for network communication or local inter-process communication.
\end{itemize}

Besides \cmd{ls -l}, two other commands are very useful for inspecting files. The \cmd{stat} command displays complete file metadata: size, hard link count, inode, last access time, modification time, and permissions in both octal and symbolic notation. The \cmd{file} command guesses a file's content type based on its contents, not its extension. This is useful when you find a file without an extension and want to know whether it is text, \term{binary}, or something else.

\subsection*{Lab 4.2: Identify Different File Types}

The goal of this lab is to distinguish file types from \cmd{ls -l} output and understand how to read file metadata thoroughly.

\begin{enumerate}
    \item Create several files of different types in the working directory:
\begin{lstlisting}[language=bash, caption={Creating a regular file and directory}]
cd ~/lab-os/chapter04-files
mkdir -p type-practice
echo "text file contents" > type-practice/notes.txt
mkdir type-practice/subfolder
ln -s type-practice/notes.txt type-practice/shortcut.txt
ls -lah type-practice/
\end{lstlisting}
    \textit{Observe:} Notice the first character on each line. \texttt{-} is for \file{notes.txt}, \texttt{d} is for \file{subfolder}, and \texttt{l} is for \file{shortcut.txt}. On the symbolic link line, you can also see where it points.

    \item Use \cmd{stat} to view complete metadata for a regular file:
\begin{lstlisting}[language=bash, caption={Viewing complete metadata with stat}]
stat type-practice/notes.txt
\end{lstlisting}

    \item Use \cmd{file} to identify content type based on file contents:
\begin{lstlisting}[language=bash, caption={Detecting file type with file}]
file type-practice/notes.txt
file type-practice/subfolder
file type-practice/shortcut.txt
file /bin/ls
\end{lstlisting}
    \textit{Notice:} See how \cmd{file} describes \file{/bin/ls} as an \textit{ELF executable}, not merely a ``file''. This helps when you encounter files without extensions.

    \item View block and character device files in \dir{/dev/}:
\begin{lstlisting}[language=bash, caption={Viewing device files in /dev}]
ls -lah /dev/tty /dev/null /dev/zero
ls -lah /dev/sd* 2>/dev/null | head -5
\end{lstlisting}
\end{enumerate}

\begin{challengebox}
Create a symbolic link named \file{home-link} in \dir{~/lab-os/chapter04-files/} that points to your home directory (\dir{~/}). Verify that the link works by running \cmd{ls -la home-link/} and comparing the result with \cmd{ls -la ~/}. Then use \cmd{stat} to check whether the inode of \file{home-link} and \dir{~/} is the same or different. Explain why the result looks that way.
\end{challengebox}

% ------------------------------------------------------------------------
\section{Hard Links and Symbolic Links}
\label{sec:links}

In Linux, one piece of data can have more than one name. This is the basic idea of a \textit{link}. Think of a book in a library. You can place one index card that points directly to the book, or create a copy of the catalog entry elsewhere that still refers to the same physical book. These two approaches correspond to a \textit{hard link} and a \textit{symbolic link}.

A \textbf{hard link} is an additional name that points to the same inode as the original file. As long as at least one hard link remains, the file data stays on disk. Because they share an inode, content changes through one name are immediately visible through the other names. Hard links can generally only be created for regular files and cannot cross different filesystems.

A \textbf{symbolic link}, or \textit{symlink}, is different. It is not an additional name for the same inode. It is a small file containing the target path. Because of that, a symlink can point to a directory, can cross filesystems, and can become a \textit{broken link} if the target is deleted or moved.

The command \cmd{ln source target} creates a hard link, while \cmd{ln -s source target} creates a symbolic link. Use \cmd{ls -li} to view inode numbers and verify whether two names share the same inode.

\subsection*{Lab 4.3: Compare Hard Links and Symbolic Links}

The goal of this lab is to understand the difference between hard links and symbolic links through inodes, file content changes, and behavior when the original target is deleted.

\begin{enumerate}
    \item Prepare one source file and create two link types:
\begin{lstlisting}[language=bash, caption={Creating a hard link and symbolic link}]
cd ~/lab-os/chapter04-files
mkdir -p link-practice
echo "initial version" > link-practice/notes.txt
ln link-practice/notes.txt link-practice/notes-hard.txt
ln -s link-practice/notes.txt link-practice/notes-soft.txt
ls -li link-practice/
\end{lstlisting}
    \textit{Observe:} \file{notes.txt} and \file{notes-hard.txt} should have the same inode number. The symlink has a different inode and shows an arrow toward its target.

    \item Change file contents through the hard link, then read from the original name:
\begin{lstlisting}[language=bash, caption={Changing contents through a hard link}]
echo "added through hard link" >> link-practice/notes-hard.txt
cat link-practice/notes.txt
\end{lstlisting}

    \item Change the original target and read contents through the symlink:
\begin{lstlisting}[language=bash, caption={Reading a target through a symbolic link}]
echo "added through original target" >> link-practice/notes.txt
cat link-practice/notes-soft.txt
\end{lstlisting}

    \item Delete the original file, then observe both links:
\begin{lstlisting}[language=bash, caption={Deleting the original target}]
rm link-practice/notes.txt
ls -li link-practice/
cat link-practice/notes-hard.txt
cat link-practice/notes-soft.txt
\end{lstlisting}
    \textit{Notice:} The hard link can still be read because it still points to the same inode. The symlink fails because its target path no longer exists.
\end{enumerate}

\begin{challengebox}
Create one file named \file{report.txt}, then create a hard link and a symbolic link for that file. Use \cmd{ls -li} to compare the inodes of all three names. After that, delete the original file and explain why one link still works while the other does not.
\end{challengebox}

% ------------------------------------------------------------------------
\section{File Search}
\label{sec:search}

Over time, a Linux system can contain thousands of files in hundreds of directories. When you need to find one configuration file or search for all large logs, browsing directories one by one is not realistic. This is where search commands matter.

The \cmd{find} command searches directly in the filesystem in real time. You can set search criteria such as file name, type, size, last modification time, and even permissions. After finding matching files, \cmd{find} can directly run a specific command on each result using the \texttt{-exec} option. This flexibility makes \cmd{find} one of the most powerful commands in Linux.

The \cmd{locate} command works differently. It searches a prebuilt index database created by \cmd{updatedb}. The result is very fast because it does not need to scan the filesystem directly. The weakness is clear: if a new file was created after the database was last updated, \cmd{locate} will not find it. Use \cmd{locate} when you need a quick answer and the file you are looking for was not just created.

Two additional commands are useful for finding programs: \cmd{which} displays the full location of a program in \var{PATH}, and \cmd{whereis} displays the location of a program's \term{binary}, manual pages, and source code.

\subsection*{Lab 4.4: Search for Files with Specific Criteria}

The goal of this lab is to use \cmd{find} with different criteria to find files efficiently and combine it with redirection techniques from Chapter~3.

\begin{enumerate}
    \item Create several test files with different names and sizes:
\begin{lstlisting}[language=bash, caption={Creating test files for search}]
cd ~/lab-os/chapter04-files
mkdir -p data-cari/logs data-cari/config data-cari/scripts
echo "isi log kecil" > data-cari/logs/app.log
echo "isi log besar" > data-cari/logs/error.log
yes "log line panjang" | head -5000 > data-cari/logs/access.log
echo "host=localhost" > data-cari/config/app.conf
echo '#!/bin/bash' > data-cari/scripts/backup.sh
\end{lstlisting}

    \item Find all \file{.log} files inside the experiment directory:
\begin{lstlisting}[language=bash, caption={Searching for files by name}]
find data-cari -type f -name "*.log"
\end{lstlisting}
    \textit{Observe:} The \texttt{-type f} option limits the search to regular files and ignores directories. The \texttt{-name} option accepts patterns with the wildcard \texttt{*}.

    \item Find files larger than 10 KB:
\begin{lstlisting}[language=bash, caption={Searching for files by size}]
find data-cari -type f -size +10k
\end{lstlisting}

    \item Find all \file{.log} files larger than 1 MB in \dir{/var/log/} and save the results to a file using redirection:
\begin{lstlisting}[language=bash, caption={Finding large logs and saving results}]
find /var/log -type f -name "*.log" -size +1M 2>/dev/null \
    > ~/lab-os/chapter04-files/log-besar.txt
cat ~/lab-os/chapter04-files/log-besar.txt
\end{lstlisting}
    \textit{Notice:} The pattern \texttt{2>/dev/null} discards ``Permission denied'' error messages that appear when \cmd{find} tries to enter inaccessible directories. You learned this redirection technique in Chapter~3.

    \item Use \cmd{find} with \texttt{-exec} to display the size of each found file:
\begin{lstlisting}[language=bash, caption={Running a command on each found file}]
find data-cari -type f -name "*.log" -exec ls -lh {} \;
\end{lstlisting}

    \item Use \cmd{which} and \cmd{whereis} to find program locations:
\begin{lstlisting}[language=bash, caption={Finding program locations}]
which ls
which bash
whereis grep
\end{lstlisting}
\end{enumerate}

\begin{challengebox}
Use \cmd{find} to search for all \file{.log} files in \dir{/var/log/} larger than 1 MB. Redirect the results to a file named \file{log-report.txt} in \dir{~/lab-os/chapter04-files/}. Make sure permission errors are ignored using \texttt{2>/dev/null}. Then use \cmd{grep}, from Chapter~2, and \cmd{wc -l} to count how many large log files were found. Write the count as a comment under your command.
\end{challengebox}

% ------------------------------------------------------------------------
\section{Compression and Archiving}
\label{sec:archive}

When you need to send dozens of files to a teammate or save a backup of a directory before making major changes, packaging all files into one bundle is the right step. This process is called archiving. When that bundle is also reduced in size, the process is called compression.

In Linux, \cmd{tar} is the main command for archiving. The name \texttt{tar} comes from \textit{tape archive}, because it was once used to store data on magnetic tape. Now \cmd{tar} is used to combine many files and directories into one \texttt{.tar} file. By default, \texttt{.tar} is not compressed. It only combines files. To compress at the same time, add the \texttt{-z} flag so \cmd{tar} automatically uses \cmd{gzip}, producing a \texttt{.tar.gz} file.

Here are the most commonly used \cmd{tar} flags:

\begin{itemize}
    \item \texttt{-c} : create a new archive
    \item \texttt{-x} : extract archive contents
    \item \texttt{-t} : list archive contents without extracting
    \item \texttt{-f file} : specify the archive filename, always required
    \item \texttt{-v} : display processed filenames, verbose mode
    \item \texttt{-z} : use gzip compression, for \texttt{.tar.gz}
    \item \texttt{-C directory} : extract to a specific directory
\end{itemize}

The \cmd{gzip} command works on a single file: it compresses one file into \texttt{.gz} and by default deletes the original file. Use the \texttt{-k} flag to keep the original file, and the \texttt{-d} flag for decompression, equivalent to \cmd{gunzip}.

The \cmd{zip} command is useful when you need to share archives with Windows or macOS users, because the \texttt{.zip} format is supported natively by all popular operating systems. Use \cmd{zip -r} to archive a directory recursively, and \cmd{unzip -l} to view archive contents without extracting them.

\subsection*{Lab 4.5: Create and Extract Archives}

The goal of this lab is to create compressed archives, inspect their contents, and extract them to a different location.

\begin{enumerate}
    \item Create a directory structure with several files to archive:
\begin{lstlisting}[language=bash, caption={Preparing a directory to archive}]
cd ~/lab-os/chapter04-files
mkdir -p proyek/src proyek/config proyek/docs
echo "# Main code" > proyek/src/main.sh
echo "PORT=8080" > proyek/config/app.conf
echo "# Documentation" > proyek/docs/readme.txt
ls -lR proyek/
\end{lstlisting}

    \item Create a compressed archive from the \dir{proyek/} directory:
\begin{lstlisting}[language=bash, caption={Creating a .tar.gz archive}]
tar -czvf proyek-backup.tar.gz proyek/
ls -lah proyek-backup.tar.gz
\end{lstlisting}
    \textit{Observe:} The \texttt{-v} flag displays every file added to the archive. Compare the compressed archive size with the original directory size.

    \item View archive contents without extracting:
\begin{lstlisting}[language=bash, caption={Viewing archive contents without extraction}]
tar -tzvf proyek-backup.tar.gz
\end{lstlisting}

    \item Extract the archive to a new directory:
\begin{lstlisting}[language=bash, caption={Extracting an archive to another directory}]
mkdir -p hasil-ekstrak
tar -xzvf proyek-backup.tar.gz -C hasil-ekstrak/
ls -lR hasil-ekstrak/
\end{lstlisting}

    \item Compress a single file using \cmd{gzip} while keeping the original file:
\begin{lstlisting}[language=bash, caption={Single-file compression with gzip}]
yes "log line" | head -10000 > long-notes.txt
ls -lah long-notes.txt
gzip -k long-notes.txt
ls -lah long-notes.txt long-notes.txt.gz
\end{lstlisting}
    \textit{Observe:} Compare the original file size with the compressed version. What percentage of space did you save?

    \item Create a ZIP archive for cross-operating-system compatibility:
\begin{lstlisting}[language=bash, caption={Creating a ZIP archive}]
zip -r proyek.zip proyek/
unzip -l proyek.zip | head -15
\end{lstlisting}
\end{enumerate}

\begin{challengebox}
Archive the entire \dir{~/lab-os/chapter04-files/} directory into a file named \file{chapter04-backup.tar.gz} stored in \dir{~/lab-os/}. After finishing, use \cmd{tar -tzvf} to verify the archive contents and make sure all subdirectories and files are included. Then use a pipe (\texttt{|}) to count how many files are inside the archive with \cmd{wc -l}. You learned this pipe technique in Chapter~3.
\end{challengebox}

% ------------------------------------------------------------------------
\section{Summary}
\label{sec:summary-filesystem}

In this chapter, you learned:

\begin{itemize}[noitemsep,topsep=2pt]
    \item \textbf{FHS hierarchy}: Linux organizes all files in one directory tree rooted at \dir{/}. Each directory has a specific function: \dir{/etc/} for configuration, \dir{/var/log/} for logs, and \dir{/home/} for user data.
    \item \textbf{File types}: The first character in \cmd{ls -l} output determines the file type. Seven major types include regular files (\texttt{-}), directories (\texttt{d}), symbolic links (\texttt{l}), and device files (\texttt{b}, \texttt{c}).
    \item \textbf{Links}: Hard links point to the same inode as the original file, while symbolic links point to a target path. The \cmd{ln} and \cmd{ln -s} commands create them.
    \item \textbf{File search}: The \cmd{find} command searches directly in the filesystem using criteria such as name, type, size, and modification time. The \cmd{locate} command searches quickly using an index database that must be updated periodically.
    \item \textbf{Compression and archiving}: The \cmd{tar} command combines many files into one archive. With the \texttt{-z} flag, \cmd{tar} also compresses using \cmd{gzip}. The \texttt{.zip} format is useful for sharing files across operating systems.
\end{itemize}

% ------------------------------------------------------------------------
\section{Lab Assignment}
\label{sec:exercises-filesystem}

\textbf{General Instructions}: Complete all exercises independently. Record important steps, save screenshots when needed, and summarize the results as personal documentation.

\subsubsection*{Lab Assignment 4.1 File System Mapping}
Run \cmd{ls -lah /} and choose eight directories from the output. For each directory, write its function in one sentence and give one real file example inside it, verified with \cmd{ls}. Include screenshot evidence for each directory you choose.


\subsubsection*{Lab Assignment 4.2 Hard Link and Symbolic Link}
Create one file \file{sumber.txt} in \dir{~/lab-os/chapter04-tugas/}. Then:
\begin{enumerate}
    \item create one hard link named \file{sumber-hard.txt};
    \item create one symbolic link named \file{sumber-soft.txt};
    \item prove with \cmd{ls -li} which names share an inode;
    \item change file contents through the hard link, then display the contents from all three names;
    \item delete \file{sumber.txt} and explain the state of the two remaining links.
\end{enumerate}
Include all commands and a summary of your observations.


\subsubsection*{Lab Assignment 4.3 Search and Archiving}
Complete the following steps in order and document the results:
\begin{enumerate}
    \item Use \cmd{find} to search for all \file{.conf} files in \dir{/etc/} smaller than 5 KB. Redirect the results to \file{config-kecil.txt} and discard permission errors.
    \item Count how many files were found using \cmd{wc -l}.
    \item Archive the entire \dir{~/lab-os/chapter04-tugas/} directory into \file{chapter04-tugas.tar.gz}.
    \item Verify the archive contents with \cmd{tar -tzvf} and count the number of files inside it using a pipe to \cmd{wc -l}.
\end{enumerate}
